% ============================================================
% THEME.TEX
% Sistema Adaptativo de Classificação Multicritério
% de Fluidos de Corte
%
% FINALIDADE
%
% Centralizar a configuração visual da apresentação Beamer.
%
% Este arquivo define apenas:
%
% - tema-base;
% - títulos dos frames;
% - rodapé;
% - numeração;
% - blocos;
% - comportamento visual geral.
%
% NÃO deve conter:
%
% - resultados científicos;
% - ranking;
% - pesos;
% - conclusões;
% - valores calculados;
% - informações institucionais não confirmadas.
%
% ============================================================


% ============================================================
% TEMA BASE
% ============================================================
%
% Enquanto não houver template oficial de apresentação
% fornecido pela professora, utilizar um tema Beamer neutro.
%
% Se posteriormente for fornecido um modelo institucional,
% ele terá prioridade sobre esta configuração.
%
% ============================================================

\usetheme{default}
\usecolortheme{default}


% ============================================================
% FONTE
% ============================================================
%
% Utilizar inicialmente a fonte sans-serif padrão do Beamer.
%
% Não impor fonte institucional enquanto não houver
% especificação oficial.
%
% ============================================================

\usefonttheme{professionalfonts}


% ============================================================
% SÍMBOLOS DE NAVEGAÇÃO
% ============================================================
%
% Removidos para manter a apresentação visualmente limpa.
%
% ============================================================

\setbeamertemplate{navigation symbols}{}


% ============================================================
% TÍTULO DOS FRAMES
% ============================================================

\setbeamertemplate{frametitle}{
    \nointerlineskip
    \begin{beamercolorbox}[
        sep=0.25cm,
        leftskip=0.2cm,
        rightskip=0.2cm,
        wd=\paperwidth
    ]{frametitle}
        \usebeamerfont{frametitle}
        \insertframetitle
    \end{beamercolorbox}
}


% ============================================================
% FONTE DO TÍTULO DOS FRAMES
% ============================================================

\setbeamerfont{frametitle}{
    size=\large,
    series=\bfseries
}


% ============================================================
% TÍTULOS E SUBTÍTULOS
% ============================================================

\setbeamerfont{title}{
    size=\LARGE,
    series=\bfseries
}

\setbeamerfont{subtitle}{
    size=\large
}

\setbeamerfont{author}{
    size=\normalsize
}

\setbeamerfont{date}{
    size=\small
}


% ============================================================
% CORPO DO TEXTO
% ============================================================

\setbeamerfont{normal text}{
    size=\normalsize
}


% ============================================================
% BLOCOS
% ============================================================
%
% Os arquivos de sections/ utilizam:
%
% block
% alertblock
%
% A configuração abaixo mantém um estilo simples e legível.
%
% ============================================================

\setbeamertemplate{blocks}[rounded][shadow=false]


% ============================================================
% TÍTULOS DOS BLOCOS
% ============================================================

\setbeamerfont{block title}{
    series=\bfseries
}

\setbeamerfont{block title alerted}{
    series=\bfseries
}


% ============================================================
% LISTAS
% ============================================================

\setbeamertemplate{itemize item}{
    \raise0.15ex\hbox{\scriptsize$\blacktriangleright$}
}

\setbeamertemplate{itemize subitem}{
    \raise0.1ex\hbox{\tiny$\blacktriangleright$}
}


% ============================================================
% ENUMERAÇÕES
% ============================================================

\setbeamertemplate{enumerate item}{
    \insertenumlabel.
}


% ============================================================
% RODAPÉ
% ============================================================
%
% Rodapé mínimo:
%
% título abreviado do projeto
% +
% número do frame
%
% Evita adicionar informações institucionais ainda não
% confirmadas.
%
% ============================================================

\setbeamertemplate{footline}{
    \leavevmode
    \hbox{
        \begin{beamercolorbox}[
            wd=.82\paperwidth,
            ht=2.7ex,
            dp=1.1ex,
            leftskip=0.3cm
        ]{author in head/foot}

            \usebeamerfont{author in head/foot}
            Classificação Multicritério de Fluidos de Corte

        \end{beamercolorbox}

        \begin{beamercolorbox}[
            wd=.18\paperwidth,
            ht=2.7ex,
            dp=1.1ex,
            rightskip=0.3cm plus1fil
        ]{date in head/foot}

            \usebeamerfont{date in head/foot}

            \insertframenumber
            \,/\,%
            \inserttotalframenumber

        \end{beamercolorbox}
    }
}


% ============================================================
% CABEÇALHO
% ============================================================
%
% Nenhum cabeçalho adicional é utilizado nesta base.
%
% ============================================================

\setbeamertemplate{headline}{}


% ============================================================
% CAPA
% ============================================================
%
% A capa é construída manualmente em:
%
% sections/00_capa.tex
%
% Portanto, não definimos um template de title page adicional.
%
% ============================================================


% ============================================================
% SUMÁRIO AUTOMÁTICO
% ============================================================
%
% Não inserir automaticamente slides de sumário entre as
% seções.
%
% Isso evita aumentar desnecessariamente o número de slides.
%
% Caso os integrantes desejem um slide de roteiro futuramente,
% ele deverá ser inserido conscientemente.
%
% ============================================================


% ============================================================
% NUMERAÇÃO DE SEÇÕES
% ============================================================
%
% As \section{} existentes em main.tex servem principalmente
% para organização interna do documento.
%
% Não criaremos divisórias automáticas nesta fase.
%
% ============================================================


% ============================================================
% MARGENS INTERNAS
% ============================================================

\setbeamersize{
    text margin left=0.6cm,
    text margin right=0.6cm
}


% ============================================================
% ESPAÇAMENTO
% ============================================================
%
% Evitar alterações globais agressivas de espaçamento.
%
% Cada frame poderá ajustar seu conteúdo localmente quando
% necessário.
%
% ============================================================


% ============================================================
% EQUAÇÕES
% ============================================================
%
% As equações metodológicas devem permanecer legíveis.
%
% Evitar:
%
% \tiny
% \scriptsize
%
% para fórmulas centrais apenas para fazê-las caber.
%
% Se necessário, dividir o conteúdo em mais de um slide.
%
% ============================================================


% ============================================================
% TABELAS
% ============================================================
%
% Não definir tamanho de fonte global para tabelas.
%
% Os integrantes deverão ajustar cada tabela conforme sua
% complexidade.
%
% Priorizar síntese em vez de redução excessiva da fonte.
%
% ============================================================


% ============================================================
% FIGURAS
% ============================================================
%
% Não aplicar molduras ou efeitos automáticos às figuras.
%
% Gráficos científicos devem manter a aparência produzida
% pela análise em R.
%
% ============================================================


% ============================================================
% CORES
% ============================================================
%
% Nenhuma paleta personalizada é definida nesta base.
%
% Motivos:
%
% 1. ainda não há template visual oficial da professora;
% 2. não queremos inventar identidade institucional;
% 3. o visual deve poder ser substituído facilmente.
%
% Se posteriormente forem fornecidas cores oficiais,
% elas deverão ser centralizadas neste arquivo.
%
% ============================================================


% ============================================================
% LOGOS
% ============================================================
%
% Nenhuma logo é carregada pelo tema.
%
% Os arquivos oficiais poderão ser inseridos na capa através de:
%
% sections/00_capa.tex
%
% usando os recursos armazenados em:
%
% assets/
%
% ============================================================


% ============================================================
% MATERIAL OFICIAL DA PROFESSORA
% ============================================================
%
% Caso posteriormente exista:
%
% - template Beamer;
% - arquivo .sty;
% - apresentação-modelo;
% - paleta;
% - fonte;
% - cabeçalho;
% - rodapé;
% - logos;
%
% a configuração oficial deverá substituir esta base genérica.
%
% A prioridade será:
%
% MATERIAL OFICIAL
% >
% THEME.TEX GENÉRICO
%
% ============================================================


% ============================================================
% ACESSIBILIDADE
% ============================================================
%
% Durante a montagem final, os integrantes deverão priorizar:
%
% - contraste suficiente;
% - fontes legíveis;
% - gráficos com rótulos claros;
% - não depender exclusivamente de cores;
% - poucos elementos por slide;
% - destaque da mensagem principal.
%
% ============================================================


% ============================================================
% RESPONSIVIDADE DA APRESENTAÇÃO
% ============================================================
%
% A apresentação foi configurada em main.tex com:
%
% aspectratio=169
%
% ou seja, formato widescreen 16:9.
%
% Caso a professora exija outro formato, a alteração deverá ser
% realizada em main.tex.
%
% ============================================================


% ============================================================
% SLIDES MUITO CHEIOS
% ============================================================
%
% Se ocorrer:
%
% Overfull \vbox
%
% ou conteúdo ultrapassando o frame, a primeira solução deve ser:
%
% - remover conteúdo redundante;
% - dividir o frame;
% - substituir texto por figura;
%
% e não simplesmente reduzir toda a fonte.
%
% ============================================================


% ============================================================
% PLACEHOLDERS
% ============================================================
%
% Durante a fase de desenvolvimento, vários frames possuem:
%
% TO_BE_FILLED
%
% Esses placeholders deverão ser substituídos ou removidos antes
% da apresentação final.
%
% ============================================================


% ============================================================
% VERSÃO FINAL
% ============================================================
%
% Antes da apresentação, conferir:
%
% - nenhum TO_BE_FILLED visível;
% - nenhum arquivo ausente;
% - nenhuma figura quebrada;
% - nenhum overflow;
% - nenhum resultado divergente do relatório;
% - numeração correta;
% - legibilidade em tela cheia;
% - logos oficiais corretas;
% - identidade visual aprovada.
%
% ============================================================


% ============================================================
% ESTADO ATUAL
% ============================================================
%
% THEME_STATUS = BASE_PREPARED
%
% OFFICIAL_PRESENTATION_THEME = NOT_PROVIDED
% CUSTOM_COLOR_PALETTE = NOT_DEFINED
% OFFICIAL_LOGOS = TO_BE_FILLED
%
% ============================================================


% ============================================================
% PRINCÍPIO FINAL
% ============================================================
%
% O tema deve apoiar a comunicação científica e nunca competir
% visualmente com os resultados.
%
% ============================================================